\DeclareMathOperator{\argmin}{arg\,min} % thin space, limits on side in displays
\DeclareMathOperator{\Cof}{Cof}
\DeclareMathOperator{\rect}{rect}
\DeclareMathOperator{\sign}{sgn}  %DESCRIPTION Sign operator
\DeclareMathOperator{\sinc}{sinc}
\DeclareMathOperator{\tr}{tr}
\DeclareMathOperator{\tri}{tri}
\DeclareMathOperator*{\mini}{\text{minimize}}
\newcommand{\christoffel}[3]{\ensuremath{\Gamma^{#1#2}_{#3}}}
\newcommand{\cross}[2]{\ensuremath{{#1} \times {#2}}}   % alternative: \wedge
\newcommand{\pdv}[2]{\frac{\partial {#1} }{\partial {#2}}} % emulate physics
\newcommand{\Curl}{\boldsymbol{\nabla}\times}
\newcommand{\Div}{\boldsymbol{\nabla}\cdot}
\newcommand{\Grad}{\boldsymbol{\nabla}}
\newcommand{\Lap}{\boldsymbol{\nabla}^2}
\newcommand{\Biharmonic}{\boldsymbol{\nabla}^4}
\newcommand{\pvint}{\mathop{\mathrm{PV}}\!\int}
\newcommand{\ray}[1]{\overset{\rightharpoonup}{#1}}


\newcommand{\abs}[1]{\left\lvert#1\right\rvert}
\newcommand{\norm}[1]{\left\lVert#1\right\rVert}
\newcommand{\del}{\partial}
\newcommand{\eps}{\varepsilon}
\newcommand{\brak}[2]{\left\langle#1,#2\right\rangle}

% other notation for spin, superscript "cross", bmatrix 

\NewDocumentCommand{\fff}{}{\mathrm{I}}





